% ============================================================
% CPET-YS Lai Lab 簡報來源文稿範本（LaTeX 版）v2.2 — 對應 skill v2.11
% 用法：複製本檔填寫後交給 cpet-yslai-pptx skill 生成簡報。
% skill 直接閱讀 .tex：\section = 章（running title）、
% \subsection = 章內一頁的頁大標、公式環境原樣轉入 plan、
% 「% [提示]」註解行 = 版面提示（與 Markdown 版語法完全相同）。
%
% v2.0 重點：新增字卡與語意項目符號，並說明排版式樣如何決定。
% v2.1 重點：補上建置前結構檢查（plan_lint）會驗的量化上限。
% v2.2 重點（ECCE 2026 回饋）：行內重點 !!片語!!、淡底字卡、流程階段箭頭、並列泳道、
%       備忘錄、備用頁、雙層標題、研討會口頭報告類型。
% ============================================================
\documentclass{article}
\usepackage{amsmath, amssymb, graphicx, CJKutf8}
\begin{document}

% ---- 簡報資訊（必填——對應 skill 步驟 0.5 使用情境確認）----
% 簡報類型：成果報告        （成果報告 / 講義 / SOP / weekly / 研討會口頭報告）
% 標題樣式：三層            （三層 / 雙層——雙層＝章名與頁名同一行，研討會官方模板建議用）
% 模板：實驗室              （實驗室 / 大會官方模板檔名——官方模板需要對應的描述檔 profile）
% 書面報告配套：是           （是→圖表章節式編號；否→流水號）
% 教學情境：否
% 語言：zh-TW               （對外口試、會議請填 en）
% 標題（英）：Your English Title Here
% 標題（中）：你的中文標題
% 報告人：姓名
% 指導教授：Y.S. Lai
% 日期：2026/07/14

% ---- 標題三層（v2.8）----
% \section    = 章 → 每頁上方置中的 running title（藍色 40pt）
% \subsection = 章內一頁 → 該頁的頁大標（navy 36pt，前置箭號）
% % [sub: 一句話] = 頁副標（灰色 16pt，選填）
%
% ---- 排版式樣（v2.8 — 不必宣告，寫什麼就排什麼）----
% 一頁由「文字／圖／字卡／表格」四種區塊依你寫下的順序垂直堆疊而成，
% 實驗室的五種式樣全部是它們的排列組合：
%   上文＋下圖          → 段落文字，接一個 % [figure:]
%   左文＋右圖          → % [layout: two_col]（左右分欄要明確指定）
%   上文＋左右兩圖      → 段落文字，接兩個 % [figure:]
%   上文＋圖＋字卡      → 段落文字，接一個 % [figure:]，接 % [card:]
%   上文＋兩圖＋兩字卡  → 段落文字，接兩個 % [figure:]，接兩個 % [card:]
% 所以不需要去記「這是式樣 4 還是式樣 5」；只有想覆寫預設判斷時才用
% % [layout: 類型]。一頁最多三個區塊，超過會警告並建議拆頁。
% 含圖的頁面，圖片會優先取得至少一半的內容區高度，字卡與條列自動讓位。

% ---- 填寫上限（v2.1）：skill 建置前會自動跑結構檢查，先遵守可少一輪修訂 ----
%   一頁的區塊數（文字／圖／字卡／表格）  <=3   ← 超過會直接中止建置，請拆頁
%   單頁 bullet 條數                      <=6
%   Outline 項數                          <=8
%   頁大標寬度（全形字算 1、半形算 0.5）  <=20  超過請把後段移到 % [sub:]
%   一頁的 [card-accent:] 強調字卡        <=1 張
%   研討會口頭報告每張字卡的條列          <=3 點（v2.2）
%   全簡報語意色（含結論橙）              <=3 種（v2.2，淡底卡與泳道的顏色都算）
%   每條公式                              必附 % [note:] 物理意義
%   取自網路或文獻的圖                    必附 source（學術誠信要求）
% 圖片路徑請用相對路徑（相對於本文稿所在目錄）；絕對路徑換機器或換工作階段
% 就會失效，skill 找不到圖檔會直接中止建置。

\section{Motivation}
\subsection{為何需要升降壓能力}
% [sub: 燃料電池與電動載具之寬電壓挑戰]
% 段落文字會被濃縮成 bullets（每段 1–2 條，一頁一訊息）
此處寫動機段落……

\begin{itemize}
\item 想強調的論點一
\item (-) 前人方法的缺陷（(-) 渲染成叉號，用於研究動機的痛點）
\item (+) 本文取得的優勢或實測結果（(+) 渲染成勾號）
\item (>) 演算法或實驗的執行步驟（(>) 渲染成箭號，用於順序性內容）
\end{itemize}
% ---- 語意項目符號（v2.8）：行首加前綴即可，全部使用 PowerPoint 內建符號 ----
%   (+) 勾號 → 研究貢獻、實驗結果驗證
%   (-) 叉號 → 前人方法缺陷、研究動機痛點
%   (>) 箭號 → 演算法執行順序、實驗步驟
% 不加前綴 = 一般箭頭項目符號。三者只用於上述語境，不要當一般裝飾。

\section{Proposed Method}
% [section: Part-I: 方法名稱一 | Part-II: 方法名稱二]
% skill 會在每個 Part 開頭放進度回訪頁（當前 Part 強調、其餘灰階）

\subsection{兩級式架構的本質限制}
% [sub: 比較基準與研究動機]
% [cards: 2]
% [card: 文獻 [11]　同規格馬達之兩級式平台 |
%   - 四開關升降壓 DC/DC + VSI
%   - (+) 低速輕載效率提升最高 23.55%]
% [card: 兩級式架構之本質限制 |
%   - (-) 兩次能量轉換損耗
%   - (-) 控制迴路相互耦合]
% [card-accent: 對策 | 單級 Y-VSI：一級內同時升降壓，內建 LC 濾波直接輸出正弦]
%
% ---- 字卡語法（v2.8）----
%   [cards: 欄數]                  宣告字卡排成幾欄（1–4，可省略）
%   [card: 標題 | 內文]            一般字卡（淺灰底）
%   [card-accent: 標題 | 內文]     強調字卡（淡橙底）——一頁最多一張
%   [card-num: 編號 | 標題 | 內文] 數字圓徽章字卡（順序本身帶訊息時才用）
%   [kpi: 86.5% | 最高效率 | BBC-DPWM @ 1800 rpm]   大數值卡
% 內文可寫段落或用「- 」條列（支援 (+)/(-)/(>) 前綴）。
% 字卡預設用箭號，與目錄／Outline 一致；要數字徽章才用 [card-num:]。
%   [card-tint: blue | 標題 | 內文]  淡底字卡（v2.2）：blue／orange／vermilion／green
%     green 是第 4 語意色，只在卡片另有 (+)/(-) 符號或文字標籤時用。
% 行內重點（v2.2）：!!片語!! → 紅色粗體；寫在 [takeaway:] 裡則是黃字放大。
%   每張卡 1–2 個、每條結論 <=3 個，重點放在結論性術語與數字。
% 流程階段箭頭（v2.2），連續 3–5 行成一組，[chevron-accent:] 為橙底強調階段：
%   % [chevron: ① 模型精修 | 細網格內插・電流限制]
%   % [chevron-accent: ② 軌跡搜尋 | MTPA／FW／MTPV 同一模型]
% 並列對照泳道（v2.2）——傳統做法 vs 本文做法，2–3 條：
%   % [lane: 傳統手調 | blue]
%   % [lane-item: figs/dyno.png | 測功機台架]
%   % [lane-item: 逐點手動調校]
%   % [lane-kpi: ≈40 h | 實驗室台架紀錄]
% 備忘錄與備用頁（v2.2）：
%   % [notes: 這頁要講的話]          寫進該頁備忘錄
%   \subsection{[backup] 頁名}       備用頁：建置時放到致謝頁之後

\subsection{方法原理}   % [layout: two_col]（文字先、公式後）
先用文字段落說明設計目的與物理直覺。

\subsection{數學模型}   % [layout: equation]
% [intro: 一句條件／目的開場句]
\begin{equation}
v_d = R_s i_d + L_d \frac{di_d}{dt} - \omega_e L_q i_q
\end{equation}
% [note: d 軸電壓方程——每條公式必附物理意義（必填）]
% [assumption: 忽略磁飽和與交叉耦合]
% 理論推導的完整中間步驟請移至文末 Remarks 節，正文只留「起點 → 假設 → 結果」

\subsection{文獻總覽與研究缺口}
\begin{tabular}{lccc}
文獻 & 架構 & 負載 & 效率驗證 \\
[4] & 單級 Y-VSI & 電阻負載 & 實測 \\
[11] & 兩級式 & SPMSM 750 W & 實測 \\
本文 & 單級 Y-VSI & 實際 SPMSM & 實測 \\
\end{tabular}
% [highlight: 本文]
% [note: 註：文獻 [10] 為架構優化分析，未進行實機驗證（註腳 11.5pt 灰字）]
% [card-accent: 研究缺口 | 尚無研究於實際 SPMSM 條件下系統性比較五種調變策略]
%
% 表格自動套用 navy 表頭＋第一欄項目欄淡藍粗體＋細框線；短內容欄自動置中。
% [highlight: 關鍵字] 把該列整列套淡橙底＋粗體（用於「本文」列或最佳結果列）。
% 要關閉第一欄樣式：% [first_col_item: false]
% 表格下方直接接 [card-accent:] 就是「表格＋結論卡」的式樣。

\section{Verification}
\subsection{各工作點最佳效率調變策略地圖}
測試條件：轉子鎖定，$i_q^* = 0.5 \cdot u(t)$ A
% [figure: figs/step_response.png | 圖 5.1 q 軸電流步階響應 | source: [4] M. Antivachis et al., IEEE JESTPE, vol. 8, 2020.]
% [kpi: 4.2 kHz | 電流迴路頻寬 | 上升時間 85 μs]
% [card: 觀察 |
%   - (+) 全部 15 個工作點皆取得最佳效率
%   - (-) 無箝位區段者效率全域偏低]
% [takeaway: 效率取決於箝位降損、THDi 取決於箝位位置，兩者相互獨立]
%
% ---- figure 語法（v2.8 擴充）----
%   [figure: 圖檔路徑 | caption | source: 出處 | takeaway: 量化結論]
% source 與 takeaway 可省略，caption 建議都寫。
% **取自網路或文獻的圖必須寫 source**（網址或 IEEE 引用格式）——這是學術誠信
% 要求，不是排版偏好；自繪圖建議標「本文實測」以與外部素材區別。
% 同頁多圖共用同一出處時，改用單獨一行 % [source: 出處] 收在該區塊下方。
% （書面配套時 caption 用章節式編號「圖 5.1」；獨立簡報用流水號「圖 1」）

\section{Conclusions}
\begin{itemize}
\item 結論一（完整句）
\item (+) 末條集中量化指標（全部 $\le 4$ 條）
\end{itemize}

\section{Remarks}
% 選填：理論推導完整步驟、冗長補充放這裡，
% 生成於 Conclusion 之後的 Remarks 章節。

% 不論簡報類型，簡報最後一律自動加致謝頁（Thank You for Your Attention，
% v2.6 起的硬規則）；要自訂文字：% [closing: 自訂感謝詞]
% [takeaway: ...] 會渲染成頁底的 navy 實心橫幅＋白字（全頁唯一的深色塊＝整頁
% 結論），與淡橙底的強調字卡層級分明。

\end{document}
